\section*{Capítulo \ref{ch:b2:magma-emplacement} --- Los magmas y su emplazamiento}

\begin{solution}{exo:b2:magma-emplacement:1}
$T = \qty{1073}{K}$. Para $w = 1$, $\ln w = 0$: $\log_{10}\eta = -3.545 +
9601/(1073 - 195.7) = -3.545 + 10.94 = 7.40$. Para $w = 4$, $\ln 4 = 1.386$:
$A = -3.545 + 1.155 = -2.390$, $B = 9601 - 3283 = 6318$, $C = 195.7 + 44.7 =
240.4$, y $\log_{10}\eta = -2.390 + 6318/832.6 = -2.390 + 7.588 = 5.20$. La
viscosidad baja en un factor $10^{7.40 - 5.20} = 10^{2.20} \approx 160$.
\end{solution}

\begin{solution}{exo:b2:magma-emplacement:2}
Riolita: $p_s = (6/0.44)^2 = \qty{186}{MPa}$, a una profundidad de $\num{186e6}/(2600\times9.81)
= \qty{7.3}{km}$. Basalto: $p_s = (0.5/0.34)^2 = \qty{2.16}{MPa}$, a
$\num{2.16e6}/(2600\times9.81) = \qty{85}{m}$. La riolita empieza a formar burbujas
en lo profundo de la corteza, y el basalto solo cerca de la superficie (si no lleva CO$_2$).
\end{solution}

\begin{solution}{exo:b2:magma-emplacement:3}
$(1150 + 50)/2 = \qty{600}{\celsius}$ y $(1150 + 400)/2 = \qty{775}{\celsius}$. El
\omterm{def:b2:magma-emplacement:bodies}{dique} en la roca fría se enfría más bruscamente en su contacto, y tiene el \omterm{def:b2:magma-emplacement:aureole}{borde de
enfriamiento rápido} más ancho.
\end{solution}

\begin{solution}{exo:b2:magma-emplacement:4}
(a) Un \omterm{def:b2:magma-emplacement:bodies}{filón capa}, paralelo a la estratificación. (b) Un \omterm{def:b2:magma-emplacement:bodies}{dique}, que corta la estratificación.
(c) Un \omterm{def:b2:magma-emplacement:bodies}{batolito}, de más de \qty{100}{km^2} y ensamblado a partir de muchas intrusiones.
(d) Un \omterm{def:b2:magma-emplacement:bodies}{lacolito}, que ha abombado su techo. (e) Un \omterm{def:b2:magma-emplacement:bodies}{dique}: el granito no tiene estratificación,
de modo que cualquier lámina en él es un \omterm{def:b2:magma-emplacement:bodies}{dique}, sea cual sea su orientación.
\end{solution}

\begin{solution}{exo:b2:magma-emplacement:5}
$(1 - 1.35\phi)^{-2.5}$: para $\phi = 0.05$, $0.9325^{-2.5} = 1.19$, es decir
\qty{119}{Pa.s}, frente a $1 + 2.5\times0.05 = 1.125$, \qty{113}{Pa.s}, con la fórmula de
Einstein (cercanos, como debe ser en el límite diluido). Para $0.20$: $0.73^{-2.5} =
2.20$, \qty{220}{Pa.s}; para $0.40$: $0.46^{-2.5} = 6.97$, \qty{697}{Pa.s}; para $0.60$:
$0.19^{-2.5} = 63.6$, unos \qty{6360}{Pa.s}. Bloqueo para $\phi = 1/1.35 = 0.74$.
\end{solution}

\begin{solution}{exo:b2:magma-emplacement:6}
$G = 300\times9.81 = \qty{2943}{Pa/m}$. (a) $\bar u = 2943\times1^2/(12\times100) =
\qty{2.45}{m/s}$; $q = \qty{2.45}{m^2/s}$; $Re = 2700\times2.45\times1/100 = 66$, laminar;
$t = 10000/2.45 = \qty{4080}{s}$, cerca de \qty{1.1}{h}. (b) $\bar u =
2943/(12\times\num{e6}) = \qty{2.45e-4}{m/s}$, $t = \qty{4.08e7}{s} = 472$ días; con
$w = \qty{10}{m}$, $\bar u$ es cien veces mayor, \qty{0.0245}{m/s}, y $t =
\qty{4.08e5}{s} = 4.7$ días. (c) El \omterm{def:b2:magma-emplacement:bodies}{dique} de \qty{1}{m} se enfría en unos
$0.5^2/\num{e-6} = \qty{2.5e5}{s}$, 2.9 días; el \omterm{def:b2:magma-emplacement:bodies}{dique} de \qty{10}{m}, en $5^2/\num{e-6} =
\qty{2.5e7}{s}$, 289 días. El basalto (\qty{1.1}{h}) y la riolita en el \omterm{def:b2:magma-emplacement:bodies}{dique}
ancho (4.7 días) llegan antes de que sus \omterm{def:b2:magma-emplacement:bodies}{diques} se congelen; la riolita en el
\omterm{def:b2:magma-emplacement:bodies}{dique} de \qty{1}{m} (472 días) no puede.
\end{solution}

\begin{solution}{exo:b2:magma-emplacement:7}
Por kilogramo de \omterm{def:b2:magma-emplacement:magma}{magma}, $n = 10/18.015 = \qty{0.555}{mol}$ de agua y $nRT =
0.555\times8.314\times1173 = \qty{5414}{J}$; el fundido restante ocupa $0.99/2400 =
\qty{4.13e-4}{m^3}$. A \qty{100}{MPa}: $V = \qty{5.41e-5}{m^3}$, fracción
$5.41/(5.41 + 41.3) = 0.12$. A \qty{10}{MPa}: $V = \qty{5.41e-4}{m^3}$, fracción
$0.57$. A \qty{0.1}{MPa}: $V = \qty{0.0541}{m^3}$, fracción $0.992$. El valor a
\qty{100}{MPa} es el menos fiable: el vapor de agua a esa presión es denso y
está lejos de ser ideal (y buena parte del agua seguiría de hecho disuelta).
\end{solution}

\begin{solution}{exo:b2:magma-emplacement:8}
(a) $2300 + 60z = 2650$ da $z = \qty{5.8}{km}$. (b) El fundido tiene
\qty{2580}{kg/m^3}: $z = 280/60 = \qty{4.7}{km}$, cerca de \qty{1.2}{km} más arriba. (c)
A $z = \qty{7.8}{km}$ la roca tiene $2300 + 60\times7.8 = \qty{2770}{kg/m^3}$; la
fuerza es $(2770 - 2650)\times9.81 = \qty{1.2}{kN}$ por metro cúbico, hacia arriba.
\end{solution}

\begin{solution}{exo:b2:magma-emplacement:9}
$v = 2\times300\times9.81\times(1000)^2/(9\times\num{e19}) = \qty{6.5e-11}{m/s}$, es decir
$\num{6.5e-11}\times\num{3.156e7} = \qty{2.1e-3}{m}$ al año, \qty{2.1}{mm/yr}. Subir
\qty{10}{km} lleva $\num{e4}/\num{2.1e-3} = \qty{4.8}{Myr}$. La escala de tiempo del enfriamiento
es $(1000)^2/\num{e-6} = \qty{e12}{s} = \qty{32}{kyr}$: el cuerpo se solidifica tras
subir unos $32000\times0.0021 \approx \qty{70}{m}$. Un \omterm{def:b2:magma-emplacement:ascent}{diapiro} magmático no puede atravesar
una corteza así; necesita un entorno más caliente y blando.
\end{solution}

\begin{solution}{exo:b2:magma-emplacement:10}
(a) $x_{\mathrm{OH}} = \sqrt{K_2x_m} = \sqrt{K_1K_2f}$, de modo que $x = K_1f + \frac12\sqrt{K_1K_2f}$.
El cociente entre el primer término y el segundo es $2\sqrt{K_1f/K_2}$: tiende a cero
cuando $f \to 0$ (domina la raíz cuadrada) y crece sin límite con $f$ (domina el
término lineal). (b) $K_1 f = \frac12\sqrt{K_1K_2 f}$ da $K_1^2f^2 =
K_1K_2f/4$, de modo que $f = K_2/(4K_1)$. (c) $f = 0.2/(\num{4e-4}) = \qty{500}{MPa}$. A
\qty{100}{MPa}: $x_m = 0.0100$, $\frac12\sqrt{\num{e-4}\times0.2\times100} =
\frac12\times0.0447 = 0.0224$, de modo que $x = 0.0324$ y la proporción de hidroxilo es
$0.0224/0.0324 = \qty{69}{\percent}$.
\end{solution}

\begin{solution}{exo:b2:magma-emplacement:11}
(a) $1000/\num{4e6} = \qty{2.5e-4}{km^3}$ al año, es decir $\qty{2.5e5}{m^3}$ al año,
o $\num{2.5e5}/\num{3.156e7} = \qty{0.0079}{m^3/s}$, unos ocho litros por segundo.
(b) El \omterm{def:b2:magma-emplacement:bodies}{dique} transporta $1000\times2\times0.8 = \qty{1600}{m^3/s}$, y aporta
$\qty{e12}{m^3}$ en \qty{6.25e8}{s}, unos 20 años: $\num{5e-6}$ de los
\qty{4}{Myr}. (c) Cada lámina contiene $20\times\num{e4}\times\num{e4} = \qty{2e9}{m^3} =
\qty{2}{km^3}$, de modo que hacen falta 500 láminas, una cada $\num{4e6}/500 = 8000$ años de media.
(d) $t = (10)^2/(4\times0.65^2\times\num{0.7e-6}) = \qty{8.5e7}{s}$, 2.7 años. Cada
lámina se solidifica miles de años antes de que llegue la siguiente: salvo que el aporte
llegue en pulsos rápidos, el \omterm{def:b2:magma-emplacement:bodies}{plutón} nunca es un único gran cuerpo de \omterm{def:b2:magma-emplacement:magma}{magma}, como sugiere la
datación de los grandes plutones.
\end{solution}

\begin{solution}{exo:b2:magma-emplacement:12}
(a) Para $x$ fija, $\ln(T - T_0) = \text{cte} - \frac12\ln t - x^2/(4\kappa t)$, cuya
derivada $-1/(2t) + x^2/(4\kappa t^2)$ se anula en $t^* = x^2/(2\kappa)$. Entonces
$e^{-x^2/(4\kappa t^*)} = e^{-1/2}$ y $\sqrt{4\pi\kappa t^*} = x\sqrt{2\pi}$, de modo que $\Delta
T_{\max} = (T_m - T_0)\,b\,e^{-1/2}/(x\sqrt{2\pi}) = (T_m - T_0)\,b/(x\sqrt{2\pi e})$. (b)
$\sqrt{2\pi e} = 4.13$: $x = 1100\times2/(100\times4.13) = \qty{5.3}{m}$, que se alcanza al cabo de
$t^* = 5.3^2/(2\times\num{0.7e-6}) = \qty{2.0e7}{s}$, unos 230 días (a solo 2.7
anchuras de \omterm{def:b2:magma-emplacement:bodies}{dique}, la aproximación de fuente plana es tosca). (c) El calentamiento decrece
como $b/x$: un \omterm{def:b2:magma-emplacement:bodies}{dique} delgado lleva poco calor, y la roca solo se calienta de forma apreciable
a unas pocas anchuras de \omterm{def:b2:magma-emplacement:bodies}{dique}.
\end{solution}

\begin{solution}{pb:b2:magma-emplacement:1}
\textbf{1.} $\log_{10}\eta = -4.55 + 5510/(1423 - 611) = -4.55 + 6.786 = 2.24$;
$\eta = \qty{172}{Pa.s}$.

\textbf{2.} $-4.55 + 5510/762 = -4.55 + 7.231 = 2.68$, $\eta = \qty{480}{Pa.s}$: un
factor $480/172 = 2.8$ en \qty{50}{K}.

\textbf{3.} $172\times(1 - 1.35\times0.20)^{-2.5} = 172\times0.73^{-2.5} = 172\times2.20 =
\qty{378}{Pa.s}$.

\textbf{4.} $\phi = 1/1.35 = 0.74$. Los datos de campo sitúan el final del flujo en masa
en una fracción de fundido de un \qtyrange{30}{50}{\percent}, es decir, un
\qtyrange{50}{70}{\percent} de cristales: antes, porque los cristales reales no son
esferas de un solo tamaño.

\textbf{5.} Basamento: $2800 - 2710 = \qty{90}{kg/m^3}$; el \omterm{def:b2:magma-emplacement:magma}{magma} es más ligero.
Sedimentos: $2450 - 2710 = \qty{-260}{kg/m^3}$; el \omterm{def:b2:magma-emplacement:magma}{magma} es más denso.

\textbf{6.} $p = 2450\times9.81\times3000 = \qty{72.1}{MPa}$.

\textbf{7.} $90\times9.81 = \qty{883}{Pa/m}$.

\textbf{8.} $\bar u = 883\times2.0^2/(12\times378) = 3532/4536 = \qty{0.78}{m/s}$.

\textbf{9.} $q = 0.78\times2.0 = \qty{1.56}{m^2/s}$; en \qty{1.0}{km}, unos
\qty{1560}{m^3/s}.

\textbf{10.} $Re = 2710\times0.78\times2.0/378 = 11$: laminar.

\textbf{11.} $12000/0.78 = \qty{15400}{s}$, unas \qty{4.3}{h}.

\textbf{12.} $(1.0)^2/\num{0.7e-6} = \qty{1.43e6}{s}$, 16.5 días, mucho más que
\qty{4.3}{h}: el \omterm{def:b2:magma-emplacement:bodies}{dique} se mantiene abierto mientras fluye el \omterm{def:b2:magma-emplacement:magma}{magma}.

\textbf{13.} A \qty{3.0}{km} el \omterm{def:b2:magma-emplacement:magma}{magma} encuentra sedimentos más ligeros que él: la
base de los sedimentos es su \omterm{def:b2:magma-emplacement:ascent}{nivel de flotabilidad neutra}. Por encima, el \omterm{def:b2:magma-emplacement:magma}{magma} sería
más denso que su entorno, de modo que se extiende lateralmente a lo largo del contacto como un
\omterm{def:b2:magma-emplacement:bodies}{filón capa}.

\textbf{14.} $p_s = (0.40/0.34)^2 = \qty{1.38}{MPa}$, que se alcanza a
$\num{1.38e6}/(2450\times9.81) = \qty{58}{m}$ bajo la superficie.

\textbf{15.} $0.34\times\sqrt{72.1} = 0.34\times8.49 = \qty{2.9}{\percent}$, mucho más que
el \qty{0.40}{\percent} presente: el \omterm{def:b2:magma-emplacement:magma}{magma} del \omterm{def:b2:magma-emplacement:bodies}{filón capa} no exsuelve agua.

\textbf{16.} El CO$_2$ es mucho menos soluble que el agua: se satura y forma burbujas
en lo profundo del ascenso, de modo que el primer gas de un basalto es sobre todo CO$_2$.

\textbf{17.} Por kilogramo, $n = 4.0/18.015 = \qty{0.222}{mol}$, $V = 0.222\times8.314\times1423/
\num{e5} = \qty{0.0263}{m^3}$; el fundido ocupa $0.996/2710 = \qty{3.68e-4}{m^3}$; la
fracción es $0.0263/(0.0263 + 0.000368) = 0.986$, un \qty{98.6}{\percent}.

\textbf{18.} $(1150 + 100)/2 = \qty{625}{\celsius}$.

\textbf{19.} Al cabo de 10 años, $2\sqrt{\kappa t} = 2\sqrt{\num{0.7e-6}\times\num{3.156e8}} =
\qty{29.7}{m}$. El punto está a $x = \qty{60}{m}$ del centro: $T = 100 +
525\,[\operatorname{erf}(-10/29.7) + \operatorname{erf}(110/29.7)] = 100 + 525\,(-0.365 +
1.000) = \qty{433}{\celsius}$.

\textbf{20.} En el centro, $T = 100 + 1050\operatorname{erf}(50/2\sqrt{\kappa t})$: al cabo de
10 años $1050\operatorname{erf}(1.68) = 1050\times0.983$, $T = \qty{1132}{\celsius}$; al cabo de
50 años, $2\sqrt{\kappa t} = \qty{66.5}{m}$, $1050\operatorname{erf}(0.752) = 1050\times0.713$,
$T = \qty{848}{\celsius}$.

\textbf{21.} En $x = \qty{130}{m}$: \qty{294}{\celsius} al cabo de 300 años,
\qty{295}{\celsius} al cabo de 363 años (el máximo) y \qty{295}{\celsius} al cabo de 420
años. Sin calor latente, la roca situada \qty{80}{m} por encima del \omterm{def:b2:magma-emplacement:bodies}{filón capa} nunca llega del todo a
\qty{300}{\celsius}.

\textbf{22.} $\num{9e8}\times1.772/(\num{2.5e6}\times1050) = 0.608$.

\textbf{23.} $\lambda = 0.64$: $e^{-0.4096}/(0.64\times1.6346) = 0.6639/1.0461 = 0.635$;
$0.65$: $0.6554/1.0673 = 0.614$; $0.66$: $0.6469/1.0886 = 0.594$. El valor $0.608$
queda entre $0.65$ y $0.66$, cerca de $0.65$: $\lambda = 0.65$.

\textbf{24.} $100 + 1050/(1 + 0.642) = 100 + 640 = \qty{740}{\celsius}$. El calor
latente liberado en el frente de solidificación mantiene el lado del \omterm{def:b2:magma-emplacement:magma}{magma} a
\qty{1150}{\celsius} más tiempo, de modo que el contacto recibe más calor.

\textbf{25.} $(50)^2/\num{0.7e-6} = \qty{3.57e9}{s}$, 113 años.

\textbf{26.} $t = (100)^2/(16\times0.65^2\times\num{0.7e-6}) = \num{e4}/\num{4.73e-6} =
\qty{2.11e9}{s}$, de modo que $\num{2.11e9}/\num{3.156e7}$: el \omterm{def:b2:magma-emplacement:bodies}{filón capa} es sólido al cabo de
\textbf{\qty[text-series-to-math=true]{67}{} años}. Es más corto que
$(b/2)^2/\kappa$ en el factor $1/(4\lambda^2) = 0.59$: esa escala de tiempo solo fija
el orden de magnitud, y las dos paredes del \omterm{def:b2:magma-emplacement:bodies}{filón capa} extraen calor hacia un encajante frío a
la vez, lo bastante deprisa para evacuar también el calor latente.
\end{solution}
